\begin{abstract}

Standard vision-language model (VLM) benchmarks evaluate whether a model can answer a question correctly given a single image, but they cannot distinguish genuine visual understanding from spurious pattern matching. We introduce \method, a \emph{counterfactual consistency} benchmark that pairs each synthetic scene with a minimal intervention---removing an object, swapping positions, or changing an attribute---and checks whether the model's answer changes \emph{if and only if} it should. \method comprises 325 programmatically generated image pairs spanning five reasoning categories (spatial, causal, compositional, counting, and occlusion) and nine intervention types, all produced with a fully deterministic PIL pipeline at zero annotation cost. We evaluate six VLMs---three proprietary (GPT-4o, Claude~3.5 Sonnet, Gemini~2.0 Flash) and three open-weight (Qwen2.5-VL 72B, Pixtral Large, Llama~3.2 11B)---using the \emph{Counterfactual Consistency Score} (\ccs), which measures the fraction of pairs where the model's answer changes only when the ground-truth answer changes. Results reveal a striking 19.5-point gap in \ccs across models (99.7\% for Gemini~2.0 Flash vs.\ 80.2\% for Llama~3.2 11B), even though all models exceed 82\% accuracy on individual images. We decompose \ccs into \emph{sensitivity} (correct change detection) and \emph{specificity} (correct invariance), uncovering distinct failure profiles: Claude~3.5 Sonnet suffers from low specificity (spurious answer changes), while Llama~3.2 11B suffers from low sensitivity (sticky answers). Notably, Qwen2.5-VL~72B achieves 97.5\% \ccs with perfect specificity, rivaling proprietary models. \method is cheap, controllable, and extensible, providing a complementary evaluation axis for VLM benchmarks. All code and data are publicly released.

\end{abstract}
